\documentclass[11pt]{article}
\usepackage[a4paper,margin=18mm]{geometry}
\usepackage{fontspec}
\setmainfont{Latin Modern Roman}
\usepackage{amsmath,amssymb,amsthm,amscd,graphicx,color}
\usepackage{xurl}
\usepackage[unicode,hidelinks]{hyperref}
\allowdisplaybreaks
\setlength\emergencystretch{3em}
\numberwithin{equation}{section}

\newtheorem{Theorem}{Theorem}[section]
\newtheorem{Lemma}[Theorem]{Lemma}
\newtheorem{Proposition}[Theorem]{Proposition}
 { \theoremstyle{definition}
\newtheorem{Remark}[Theorem]{Remark} }

\newcommand{\rme}{\mathrm{e}}
\newcommand{\const}{\mathrm{const}}
\DeclareMathOperator{\Complex}{\mathbb{C}}
\DeclareMathOperator{\Integer}{\mathbb{Z}}
\DeclareMathOperator{\rank}{rank}
\DeclareMathOperator{\CP}{\mathbb{C}P}

\DeclareMathOperator{\res}{res}
\DeclareMathOperator{\id}{id}
\DeclareMathOperator{\Ibb}{\mathbb{I}}
\DeclareMathOperator{\wgt}{wgt}
\newcommand{\rmd}{\mathrm{d}}
\newcommand{\rmi}{\mathrm{i}}
\newcommand{\Arf}{\mathrm{Arf}}
\newcommand{\Jac}{\mathrm{Jac}}


\makeatletter\newlabel{IntArbP}{{5.1}{}{Original source locator}{}{}}\makeatother
\makeatletter\newlabel{A:BLRs}{{D}{}{Original source locator}{}{}}\makeatother
\makeatletter\newlabel{A:Diffs}{{A}{}{Original source locator}{}{}}\makeatother
\makeatletter\newlabel{A:singPart}{{B}{}{Original source locator}{}{}}\makeatother
\begin{document}
\begin{center}\Large For a smooth complex genus-three (3,4) curve in form (2.4), the second-kind cohomology basis (3.2) with the primitive-sigma normalization has regularization vector (0,-lambda2/3,-lambda5/6) at infinity\end{center}
\noindent\textbf{Stable ID:} 1784\par
\noindent\textbf{Authors:} Julia Bernatska; Dmitry Leykin\par
\noindent\textbf{Original work:} On Regularization of Second Kind Integrals\par
\noindent\textbf{Version:} Official SIGMA 14 (2018) article 074, DOI 10.3842/SIGMA.2018.074; journal PDF and native TeX acquired 2026-10-01, exact bytes pinned by SHA256\par
\noindent\textbf{Selected task scope:} Full genus-three Theorem4.1 with its complete alternative Proof2. Work with a smooth nondegenerate complex curve y\textasciicircum{}3+\allowbreak{}x\textasciicircum{}4+\allowbreak{}lambda2*\allowbreak{}y*\allowbreak{}x\textasciicircum{}2+\allowbreak{}lambda5*\allowbreak{}y*\allowbreak{}x+\allowbreak{}lambda6*\allowbreak{}x\textasciicircum{}2+\allowbreak{}lambda8*\allowbreak{}y+\allowbreak{}lambda9*\allowbreak{}x+\allowbreak{}lambda12=\allowbreak{}0, the exact first/second-kind differential basis3.2, local parameter x=\allowbreak{}-xi\textasciicircum{}-3 at infinity and the multivariate sigma/primitive normalization2.3/3.1. Then c=\allowbreak{}(0,-lambda2/3,-lambda5/6). No arbitrary cohomology basis, alternate sigma normalization, singular member, other (n,s)-curve variant or addition-formula neighbor is asserted.\par
\noindent\textbf{Difficulty:} H2: The additive regularization vector is fixed by compatibility with the Abelian-function field and the specified multivalued primitive, not by removing the pole alone. The author derives residue identities and rational-function elimination relations on a nonspecial divisor, then uses a bilinear sigma identity to determine every component uniquely. This supplies the complete geometric and Abelian-function support for the final local coefficient calculation.\par
\noindent\textbf{Editorial boundary note (not author text):} The complete original genus-three Theorem 4.1 is retained with the full alternative Proof 2, including every residue/rational-function elimination relation of its new Lemma 4.5 and the final bilinear sigma/primitive Laurent expansion determining all three constants. The exact nondegenerate curve, local x=\allowbreak{}-xi\textasciicircum{}-3 coordinate, selected first/second-kind basis and primitive/sigma normalization are included. The first alternative proof remains uncounted context for the author P-function definition and normalization, rather than an assumed missing proof. Other curves, arbitrary-pole regularizations and addition-formula neighbors are unselected; no coefficient formula or proof is generated.\par
\noindent\textbf{Source:} \url{https://sigma-journal.com/2018/074/}\quad DOI: \url{https://doi.org/10.3842/SIGMA.2018.074}\par
\noindent\textbf{License and attribution:} Copyright retained by the named authors. Published by SIGMA under CC BY 4.0: \url{https://creativecommons.org/licenses/by/4.0/}. Source: \url{https://sigma-journal.com/about.html}. Adaptation consists of standalone excerpt formatting, cover and numbering restoration; original author language and mathematical text are retained.\par
\noindent\textbf{Editorial symbol notice (not author text):} The alternate Proof1 native line280 prints -log(partial\_u1\textasciicircum{}2 P)=\allowbreak{}3 log(xi)+\allowbreak{}..., whereas Lemma4.2 identifies minus the derivative with a primitive having a cubic zero. That literal logarithm sign is inconsistent and is not silently repaired. The packet uses only the fully supplied independent Proof2 on physical pages8-10; no inference from the defective displayed line is needed. Author source pages remain unchanged.\par
\noindent\textbf{External original locator IntArbP:} Original Section 5.1 extends regularization to a pole at an arbitrary point, outside the selected infinity-basepoint task.\par
\noindent\textbf{External original locator A:BLRs:} Original Appendix D contains the separate genus-four bilinear expansion, mentioned only in the overview.\par
\noindent\textbf{External original locator A:Diffs:} Original Appendix A lists differential bases for the other unselected curves.\par
\noindent\textbf{External original locator A:singPart:} Original Appendix B lists singular parts for the other unselected curves.\par
\medskip\hrule\medskip\noindent\textbf{Author text begins below.}\par
\setcounter{section}{0}
\setcounter{subsection}{0}
\setcounter{subsubsection}{0}
\setcounter{Theorem}{0}
\setcounter{equation}{0}
% Original source lines 64--449
\section{Introduction}

In this paper we consider second kind integrals on algebraic curves, which are second kind $0$-forms $r(x,y;[\gamma])$ obtained from second kind differential forms $\rmd r$
\begin{gather}\label{SecKIntA}
 r(x,y;[\gamma]) = \int_{[\gamma]}^{(x,y)} \rmd r(\tilde{x},\tilde{y}),
\end{gather}
where $[\gamma]$ denotes a class of homotopically equivalent paths {from a fixed basepoint to $(x,y)$, by $(\tilde{x},\tilde{y})$ we denote a point which serves as a dummy parameter of integration}. Definitions of second kind integrals were discussed in~\cite{Hodge}, where the theory of second kind differentials was developed from the theory of stacks adapted for use in the theory of functions of several complex variables. Here we define a second kind differential form as a locally exact meromorphic form, whose all singularities are poles of order greater than one.

In what follows we deal with a particular class of curves called $(n,s)$-curves, possessing a~Weierstrass point at infinity, which is the branch point where all sheets of the curve come together. We assign the infinity to the basepoint of Abel's map on the curve. An accurate definition of the $(n,s)$-curve is given below in Section~\ref{ss:nsCurve}. Throughout the paper we consider second kind differential forms with the only singularity at the infinity, except Section~\ref{IntArbP} where the singularity is located at an arbitrary point. We call the antiderivative of a second kind $1$-form $\rmd r$ with the basepoint at the infinity a second kind integral $r$, namely,
\begin{gather}\label{SecKInt}
 r(x,y) = \int_{\infty}^{(x,y)} \rmd r(\tilde{x},\tilde{y}),
\end{gather}
and define the latter on the fundamental domain of a genus $g$ $(n,s)$-curve. Suppose the curve has a~homology basis of $2g$ cycles $\{\mathfrak{a}_i, \mathfrak{b}_i \,|\, i=1,\dots,g\}$, all disjoint except for each pair~$\mathfrak{a}_i$ and~$\mathfrak{b}_i$ intersecting at one point. The fundamental domain is the one-connected domain obtained by cutting the curve along all~$2g$ cycles of the homology basis and $g$ paths connecting the basepoint with $g$ intersection points. Evidently, the integral~\eqref{SecKInt} requires regularization at the basepoint, which is the only singularity point of the integrand function.

The idea of regularization is adopted from the definition of Weierstrass zeta-function. On the Weierstrass normal form of elliptic curve $0=f(x,y)=-y^2 + 4x^3 - g_2 x - g_3$, which is the simplest $(n,s)$-curve, the second kind integral is defined by
\begin{gather*}
 r(x,y) = \int_{\infty}^{(x,y)} \frac{\tilde{x}\, \rmd \tilde{x}}{\partial_{\tilde{y}} f}.
\end{gather*}
Here and throughout the text we denote $\partial f/\partial y$ by $\partial_y f$. Let
\begin{gather*} (x,y)=\big(\wp(u),\wp'(u)\big)= \big(u^{-2},{-}2u^{-3}\big)+O(1)\end{gather*} be a parametrization in the vicinity of infinity, namely $(x,y)\to\infty$ as $u\to 0$, where $u$ serves as a local parameter. With the standard Abel's map $\mathcal{A}(x,y) = \int_{\infty}^{(x,y)} (\partial_{\tilde{y}} f)^{-1}\rmd \tilde{x}$ the second kind integral transforms into zeta function $\zeta \big(\mathcal{A}(x,y)\big) = r(x,y)$. The textbook definition of Weierstrass zeta-function, see, e.g., \cite[Chapter~XX, Section~20.4]{Whittaker}, serves as a regularization of the second kind integral
\begin{gather*}
 \zeta(u) = \frac{1}{u} - \int_{0}^{u} \bigg(\wp(\tilde{u}) - \frac{1}{\tilde{u}^2}\bigg) \rmd \tilde{u}.
\end{gather*}

In the present paper we extend this regularization to a wider collection of curves. In this connection, a constant arises to be added to the regularized second kind integral, we call it a regularization constant. As shown in \cite[Section~193]{BakerAF}, where the notion of a regularization constant arose, these constants vanish in the case of hyperelliptic curves. This can explain why it was not revealed before. In the non-hyperelliptic case, the regularization constant does not vanish, and plays an important role when the second kind integral is used to obtain relations between Abelian functions defined on the Jacobi variety of the curve. In particular, for this purpose we employ the primitive function
\begin{gather}\label{PsiDef0}
 \psi(x,y) = \exp\left\{-\int_{\infty}^{(x,y)} r(\tilde{x},\tilde{y})^{t} \rmd u(\tilde{x},\tilde{y}) \right\}
\end{gather}
introduced in \cite{BL2005}, where $\{\rmd u(x,y),\rmd r(x,y)\}$ form a special cohomology basis, and $r(x,y)$ is the antiderivative~\eqref{SecKInt} of the $\rmd r(x,y)$. In general, the primitive function $\psi(x,y;[\gamma])$ also depends on a~path $\gamma$ from the basepoint to a point~$(x,y)$, and the second kind integral is defined by~\eqref{SecKIntA}. As above, $[\gamma]$ denotes a class of homotopically equivalent paths.

The paper is organized as follows. In Section~\ref{s:Prel} we briefly recall the notions of $(n,s)$-curve and primitive function, we explain how to construct the special cohomology basis, and specify the curves under consideration. Section~\ref{s:RegInt} is devoted to the idea of regularization of the second kind integral~\eqref{SecKInt}, also the special cohomology bases on the curves under consideration are specified.

In Section \ref{s:Reg34} we show how to produce relations between Abelian functions defined on the Jacobian of $\mathcal{V}_{(3,4)}$ with the help of equality \eqref{BLRC34}, which involves the primitive function~$\psi$. Si\-mi\-lar computations in the case of $(3,5)$-curve are shown in Appendix~\ref{A:BLRs}. We emphasize that this method gives consistent relations between Abelian functions only with a correct choice of regularization constant in the definition of $r(x,y)$. Therefore, we could say that the regularization constant is unique, and consistent with the structure of Abelian function field.

The following proposition summarizes results obtained in the paper. Each regularization constant $c_{(n,s)}$ relates to the second kind integral on $(n,s)$-curve with the special cohomology basis.
\begin{Proposition}
In some non-hyperelliptic cases regularization constants are the following
 \begin{gather*}
 c_{(3,4)}(\lambda) = \big(0, -\tfrac{1}{3}\lambda_2, {-}\tfrac{1}{6}\lambda_5 \big)^t,\\
 c_{(3,5)}(\lambda) = \big({-}\tfrac{1}{2}\lambda_1, -\tfrac{4}{15}\lambda_1^2, -\tfrac{1}{3}\lambda_4, {-}\tfrac{1}{6}\lambda_7 \big)^t,\\
 c_{(3,7)}(\lambda) = \big(0,{-}\tfrac{2}{3}\lambda_2,{-}\tfrac{2}{7}\lambda_2^2,{-}\tfrac{1}{2}\lambda_5,
 {-}\tfrac{1}{3}\lambda_8,{-}\tfrac{1}{6}\lambda_{11}\big)^t,\\
 c_{(4,5)}(\lambda) = \big(0, \tfrac{7}{45}\lambda_2, -\tfrac{3}{4}\lambda_3, \tfrac{3}{10}(\lambda_6-\lambda_3^2),
 -\tfrac{1}{2}\lambda_7, {-}\tfrac{1}{4}\lambda_{11} \big)^t.
\end{gather*}
\end{Proposition}
We do not have a general formula for $c_{(n,s)}$, however we have developed a technique of finding~$c_{(n,s)}$ for virtually any pair~$(n,s)$. Below we propose two methods of computing $c_{(n,s)}$. We demonstrate both of them by carrying out calculation for the case of $(3,4)$-curve. Other cases are much more cumbersome.

\section{Preliminaries}\label{s:Prel}

\subsection[$(n,s)$-curve]{$\boldsymbol{(n,s)}$-curve}\label{ss:nsCurve}
As mentioned above, in the paper we deal with $(n,s)$-curves \cite{BEL1999}.
They are defined for co-prime integers $n$ and $s$
as \begin{gather*}\mathcal{V}_{(n,s)}=\big\{(x,y)\in \Complex^2 \,|\, f_{(n,s)}(x,y) =0\big\},\end{gather*} where
\begin{gather}\label{nsCurve}
 f_{(n,s)}(x,y) = y^n + x^s + \sum_{i=0, j=0}^{s-2,n-2} \lambda_{ns-in- js} x^i y^j,
\end{gather}
$\lambda_{j}\in \Complex$, $\lambda_{k}=0$ whenever $k\leqslant 0$. Nondegenerate $(n,s)$-curve has genus $g = \frac{1}{2}(n - 1)(s - 1)$, {this is the maximal genus of the family of curves of the form~\eqref{nsCurve}}, see~\cite{BEL1999}. The polynomial $f_{(n,s)}$ is homogeneous with respect to Sato weight defined by $\wgt x = n$, $\wgt y=s$, thus $\wgt \lambda_j=j$. In the vicinity of $(x,y)=\infty$ we introduce a parameterization on $\mathcal{V}_{(n,s)}$ by the following formulas
\begin{gather*}%\label{xyParam}
x(\xi)=-\xi^{-n}, \qquad y(\xi)=-\xi^{-s}\left((-1)^s+ \sum_{i=1}^{\infty} \mu_i(\lambda) \xi^i\right), \qquad f_{(n,s)}\big(x(\xi),y(\xi)\big)=0.
\end{gather*}
Note that $\wgt \xi = -1$, and $\mu_i(\lambda)$ are polynomials in $\lambda$ with rational coefficients, $\wgt \mu_i(\lambda) = i$.

\subsection{Basis of 1-forms}\label{ss:CoHomolBasis}
Fix the basis of holomorphic $1$-forms $\rmd u = \big(x^i y^j (\partial_y f_{(n,s)})^{-1}\rmd x\big)$, $i\geqslant 0$, $j\geqslant 0$ and $i n+j s\leqslant 2g-2$ ordered by descending Sato weight. Note that $\wgt \rmd u = ({-}w_1,\dots,{-}w_g)^t$, where $\{w_1,\dots,w_g\}$ is Weierstrass sequence obtained from $n$ and $s$ with the help of generating function $t^{w_1}+\cdots +t^{w_g}= (1-t)^{-1}- (1-t^{ns}) (1-t^n)^{-1}(1-t^s)^{-1}$. Further consider meromorphic $1$-forms with poles only at infinity. Up to globally exact forms any such $1$-form can be represented as a linear combination of $2g$ differentials $\big(x^i y^j (\partial_y f_{(n,s)})^{-1} \rmd x\big)$ with $0\leqslant i \leqslant s-2$ and $0\leqslant j \leqslant n-2$, which are holomorphic on $\mathcal{V}_{(n,s)}\backslash \infty$, and the condition $i n+j s\geqslant 2g$ singles out differentials that actually have a pole at infinity. One can form a vector $\rmd r$ of $\wgt \rmd r(\xi) = (w_1,\dots,w_g)^t$ which is subject to condition
\begin{gather}\label{DiffAjst}
\res_{\xi=0} \left(\int_0^{\xi} \rmd u(\tilde{\xi})\right) \rmd r(\xi)^t = 1_g.
\end{gather}
This condition completely determines the {principle part of $\rmd r(\xi)$.} Though the holomorphic part of $\rmd r(\xi)$ is inessential in what follows, it is {worth} to note that it can be chosen so that $\rmd r(x,y)$ and $\rmd u(x,y)$ form an associated system, see \cite[Section~138]{BakerAF} and \cite[p.~131]{EEL2000}. Below in actual calculations we use the {first $\mathcal{A}(x,y)$ and second kind integrals $\mathcal{A}^\ast(x,y)$ obtained from the chosen basis of $1$-forms}, namely,
\begin{gather*}\mathcal{A}(x,y) = \int_{\infty}^{(x,y)} \rmd u(\tilde{x},\tilde{y}),\qquad \mathcal{A}^\ast(x,y) = \int_{\infty}^{(x,y)} \rmd r(\tilde{x},\tilde{y}).\end{gather*}
So $\mathcal{A}\colon \mathcal{V} \mapsto \Jac(\mathcal{V})$ denotes the standard Abel's map. The meromorphic map $\mathcal{A}^\ast$ has a pole at infinity, and requires regularization. The regularization constants~$c_{(n,s)}$ computed in this paper relate to these particular second kind integrals.

\subsection{Primitive function}
The primitive function introduced by \eqref{PsiDef0} is employed by the both methods of computing the regularization constant. More accurately, we define it as
\begin{gather}\label{PsiDef}
 \psi(x,y) = \exp\left\{-\int_{\infty}^{(x,y)} \mathcal{A}^\ast(\tilde{x},\tilde{y})^{t} \rmd \mathcal{A}(\tilde{x},\tilde{y}) \right\}.
\end{gather}
In the definition an antiderivative $\mathcal{A}^\ast(x,y)$ of the second kind differential is supposed to contain the regularization constant when parametrized. In terms of local parameter $\xi$ the following representation of~\eqref{PsiDef} can be obtained
\begin{gather*}%\label{PsiRegDef}
\psi(\xi) = \xi^g \exp \left\{-\int_0^{\xi} \big(\big(\mathcal{A}^\ast(\xi)\big)^t \rmd \mathcal{A}(\xi) + g \xi^{-1}\rmd\xi \big)\right\},
\end{gather*}
where due to condition \eqref{DiffAjst} the integrand is a holomorphic function of~$\xi$. Evidently, the primitive function is entire with a~zero at $\xi=0$ of order~$g$, where~$g$ denotes the genus of a curve.

The primitive function $\psi(\xi)$ coincides with a certain derivative of sigma function $\sigma(u)$ at $u= \mathcal{A}(\xi)$ up to a constant factor, see Remark~\ref{R:SigmaDer}, and also with a certain modification of the prime form arisen from~\cite{N2010}. Following~\cite{BL2004}, we define the sigma function of $g$ variables, which we call a~multivariate sigma function, as a~solution of a~system of heat equations with a~Schur--Weierstrass polynomial as an initial condition. Here we apply the approach to constructing multivariate sigma functions on $(n,s)$-curves after~\cite{BL2004,BL2008}. The most complete survey of the theory of multivariate sigma is given in~\cite{BEL2012}.

\subsection{Curves under consideration}
The non-hyperelliptic plane algebraic curves under consideration are unfoldings of simple Pham singularities $x^s + y^n$, that is $(n,s)=(3,4)$ or $(3,5)$. Denote by $\mathcal{V}_{(3,4)}$ a genus $3$ curve defined by the equation $f(x,y)=0$, where the polynomial~$f$ is given by
\begin{gather}\label{C34eq}
f(x,y) = y^3 + x^4 + \lambda_2 y x^2 + \lambda_5 y x + \lambda_6 x^2 + \lambda_8 y + \lambda_9 x + \lambda_{12}.
\end{gather}
Denote by $\mathcal{V}_{(3,5)}$ a genus $4$ curve defined by the equation $f(x,y)=0$, where
$f$ is given by
\begin{gather}\label{C35eq}
 f(x,y) = y^3 + x^5 + \lambda_1 y x^3 + \lambda_4 y x^2 + \lambda_6 x^3 + \lambda_7 y x + \lambda_9 x^2 + \lambda_{10} y + \lambda_{12} x + \lambda_{15}.
\end{gather}
We also consider curves $\mathcal{V}_{(3,7)}$ and $\mathcal{V}_{(4,5)}$ of genus 6 with polynomials $f$ given by
\begin{gather*}
f(x,y) = y^3 + x^7 + \lambda_2 y x^4 + \lambda_5 y x^3 + \lambda_6 x^5 + \lambda_8 y x^2 + \lambda_9 x^4 + \lambda_{11} y x + \lambda_{12} x^3 \\
\hphantom{f(x,y) =}{} + \lambda_{14} y + \lambda_{15} x^2 + \lambda_{18} x + \lambda_{21},\\
f(x,y) = y^4 + x^5 + \lambda_2 y^2 x^2 + \lambda_3 y x^3 +\lambda_6 y^2 x+\lambda_7 y x^2 + \lambda_8 x^3 +\lambda_{10} y^2 +\lambda_{11} y x \\
\hphantom{f(x,y) =}{} + \lambda_{12} x^2 + \lambda_{15} y + \lambda_{16} x + \lambda_{20},
\end{gather*}
respectively.

\section{Regularization of second kind integral in general}\label{s:RegInt}

Now we define the regularization more accurately. Consider a second kind integral $r(x,y)$ in the vicinity of its pole. Let $\xi$ be a local coordinate, and $\rmd r(\xi)=r'(\xi)\rmd \xi$, where $r(\xi)$ is $0$-form. Then~$r(\xi)$ is decomposed into singular and regular parts:
\begin{gather*}%\label{Dif2Decomp}
 r(\xi) = r_{\textup{sing}}(\xi^{-1}) + r_{\textup{reg}}(\xi) + c,\qquad r_{\textup{sing}}(0)=0,\qquad r_{\textup{reg}}(0)=0.
\end{gather*}
Essentially, this decomposition is a textbook regularization, and $c$ denotes the \emph{regularization constant}. We define the \emph{regularized second kind integral} as follows
\begin{gather}\label{IntReg}
 \mathcal{A}^\ast(\xi)= r_{\textup{sing}} (\xi)
+ c(\lambda) + \int_0^{\xi} \bigl(\rmd r(\tilde{\xi})- \rmd r_{\textup{sing}} (\tilde{\xi}) \bigr).
\end{gather}
The regularization constant $c(\lambda)$ {is a vector with $g$ components, which are} polynomials in $\lambda$ with rational coefficients, and $\wgt c(\lambda) = \wgt r(\xi) = (w_1,\dots,w_g)^t$. Note that representation~\eqref{IntReg} is essentially connected to parameterization in the vicinity of infinity, and holds true within the fundamental domain. The regularization constant should not be confused with a constant of integration. The latter vanishes in order to keep $\psi$ satisfying the functional equation~\eqref{PsiFEq}. The regularization constant appears only when the second kind integral has the pole at infinity and is parameterized in the vicinity of this pole.

If one needs to define the regularization constant, we say that the correct choice of this constant makes the primitive function $\psi$ at $\xi$ equal, up to a rational factor, to the first non-vanishing derivative of sigma-function on the Abel's image of~$\xi$. We also recall {the method} of producing relations between Abelian functions, which involves the primitive function~$\psi$. The correct choice of the regularization constant is necessary for obtaining consistent relations. In what follows, we use both properties of the primitive function $\psi$ to compute the regularization constants announced in Introduction.

In the hyperelliptic case the regularization constants vanish due to the hyperelliptic involution $\iota(x,y) = (x,-y)$, that is $ = \mathcal{A}(\iota(x,y)) = -\mathcal{A}(x,y)$, or after parameterization $\mathcal{A}(-\xi) = -\mathcal{A}(\xi)$, and the fact that zeta funtions, and also second kind integrals, are odd, for more details see \cite[Section~193]{BakerAF}. Therefore, $0=r(\xi)+r(-\xi)=2c_{(2,2g+1)}(\lambda)$,

\subsection[Basis differentials and integrals on $\mathcal{V}_{(3,4)}$]{Basis differentials and integrals on $\boldsymbol{\mathcal{V}_{(3,4)}}$}
On the curve $\mathcal{V}_{(3,4)}$ punctured at infinity we fix a basis of holomorphic differentials
\begin{subequations}\label{Diffs34}
\begin{gather}
\rmd u(x,y,\lambda) = \begin{pmatrix}
 y \\ x \\ 1
 \end{pmatrix} \frac{\rmd x}{\partial_y f},\label{Diff1s34} \\
 \rmd r(x,y,\lambda) = \begin{pmatrix}
 -x^2 \\ -2 y x \\ -5 y x^2 + \frac{2}{3} \lambda_2^2 x^2 + \frac{2}{3} \lambda_5 \lambda_2 x - \lambda_6 y
 \end{pmatrix}\frac{\rmd x}{\partial_y f}. \label{Diff2s34}
\end{gather}
\end{subequations}
Introduce a local coordinate $\xi$ in the vicinity of infinity. Then we have a parameterization
\begin{gather}\label{xyParamC34}
 x(\xi) = -\xi^{-3},\qquad y(\xi) = \xi^{-4}\big({-}1 + \tfrac{1}{3}\lambda_2 \xi^2 - \tfrac{1}{3} \lambda_5 \xi^5 + O\big(\xi^6\big)\big).
\end{gather}
In terms of the local coordinate the basis holomorphic differentials acquire the form
\begin{subequations}
 \begin{gather}\label{C34D1k}
 \rmd u(\xi) =
 \begin{pmatrix}
 -1 + \frac{1}{9} \lambda_2^2 \xi^4 + O\big(\xi^6\big)\vspace{1mm}\\
 -\xi - \frac{1}{3} \lambda_2 \xi^3 + \frac{1}{3} \lambda_5 \xi^6 + O\big(\xi^7\big)\vspace{1mm}\\
 \xi^4 + \frac{1}{3} \lambda_2 \xi^6 - \frac{1}{3} \lambda_5 \xi^{9} + O\big(\xi^{10}\big)
 \end{pmatrix} \rmd \xi,\\
\label{C34D2k}
 \rmd r(\xi) =
 \begin{pmatrix} -\xi^{-2} - \frac{1}{3} \lambda_2 + \frac{1}{3} \lambda_5\xi^3 + O\big(\xi^{4}\big) \vspace{1mm}\\
 - 2 \xi^{-3} + \frac{2}{9}\lambda_2^2 \xi + O(\xi^3)\vspace{1mm}\\
 5 \xi^{-6} + \frac{1}{9} \lambda_2^2 \xi^{-2} + O(1)\end{pmatrix} \rmd \xi.
\end{gather}
\end{subequations}
The first kind integral $\mathcal{A}(\xi)=\int_0^{\xi} \rmd u(\xi)$ is a well defined holomorphic function {of local parameter~$\xi$}. On the curve $\mathcal{V}_{(3,4)}$ the regularized second kind integral $\mathcal{A}^\ast(\xi)$, defined by~\eqref{IntReg}, has the following singular part
\begin{gather*}%\label{Rsing34}
 r_{\textup{sing}} (\xi) =\begin{pmatrix} \xi^{-1} \\ \xi^{-2}\\ - \xi^{-5} - \frac{1}{9}\lambda_2^2 \xi^{-1}\end{pmatrix} ,\qquad
 \rmd r_{\textup{sing}} (\xi) = r'_{\textup{sing}} (\xi) \rmd \xi = \begin{pmatrix} - \xi^{-2} \\ - 2 \xi^{-3}\\
 5 \xi^{-6} + \frac{1}{9} \lambda_2^2 \xi^{-2}\end{pmatrix} \rmd \xi.
\end{gather*}
The integral on the right hand side of \eqref{IntReg} is regular, and $\rmd \mathcal{A}^\ast(\xi)/\rmd \xi = \rmd r(\xi)/\rmd \xi$.

Basis differentials of the first and second kinds defined on $\mathcal{V}_{(3,5)}$, $\mathcal{V}_{(3,7)}$, and $\mathcal{V}_{(4,5)}$ are given in Appendix~\ref{A:Diffs}, and singular parts of the second kind integrals in terms of a local coordinate in the vicinity of infinity can be found in Appendix~\ref{A:singPart}.

\section[Regularization of second kind integral on $(3,4)$-curve]{Regularization of second kind integral on $\boldsymbol{(3,4)}$-curve}\label{s:Reg34}
\begin{Theorem}\label{T:Reg34}
In the definition of regularized second kind integral \eqref{IntReg} on $\mathcal{V}_{(3,4)}$, with the basis first and second kind differentials given by \eqref{Diffs34}, the regularization constant is equal to
 \begin{gather}\label{C34const}
 c(\lambda) = \bigg(0, -\frac{\lambda_2}{3}, {-}\frac{\lambda_5}{6} \bigg)^t.
\end{gather}
\end{Theorem}
We propose two methods of proof for this type of assertions. One of them relies on vanishing properties of sigma-functions. The other one uses bilinear Hirota type differential equations satisfied by Abelian functions.
\begin{proof}[Proof~1] Consider the function $P(\xi,u)=\sigma\bigl(\mathcal{A}(\xi)-u\bigr)$. By Riemann vanishing theorem, $P(\xi,u)$ has at most $g=3$ zeros as a function of $u$. Let $u=\mathcal{A}(\xi_1)+\mathcal{A}(\xi_2)+\mathcal{A}(\xi_3)$ then \begin{gather*}
P(\xi,u) = (\xi-\xi_1)(\xi-\xi_2)(\xi-\xi_3)\sum_{k=0} \alpha_k (\xi_1,\xi_2,\xi_3,\lambda)\xi^k,
\end{gather*}
where $\alpha_k$ are entire series in $\lambda$ with symmetric polynomials in $\xi_1$, $\xi_2$, $\xi_3$ as coefficients. In particular, $\alpha_0 (\xi_1,\xi_2,\xi_3,0) = -\big(\xi_1^2+\xi_2^2+\xi_3^2 +\xi_1 \xi_2 + \xi_1 \xi_3 + \xi_2\xi_3\big)$, and $\alpha_k(\xi_1,\xi_2,\xi_3,0)=0$ at $k>0$, further $\alpha_k(0,0,0,\lambda)=0$ at $k\geqslant 0$. According to Riemann vanishing theorem, $P(\xi,u)$ vanishes identically if $\xi_1+\xi_2+\xi_3=0$ and $\xi_1\xi_2+ \xi_1\xi_3 + \xi_2 \xi_3 = 0$, that is $u\equiv 0$, since the points $\xi_1$, $\xi_2$, $\xi_3$ correspond to zeros of the function $x-a$ on the curve $\mathcal{V}_{(3,4)}$ for some $a\in \Complex$. Otherwise, $P(\xi,u)$ has exactly $g=3$ zeros (with multiplicities).

On the other hand, inverting Abel's map we come to the determinant
\begin{gather*}
 W = \begin{vmatrix}
 y(\xi_1) & x(\xi_1) & 1 \\ y(\xi_2) & x(\xi_2) & 1\\ y(\xi_3) & x(\xi_3) & 1
 \end{vmatrix} = \frac{(\xi_1-\xi_2)(\xi_2-\xi_3)(\xi_3-\xi_1)}{\xi_1^4 \xi_2^4 \xi_3^4}
 \beta (\xi_1,\xi_2,\xi_3,\lambda).
\end{gather*}
Similarly to $\alpha_k$, series $\beta$ is entire in $\lambda$, and vanishes identically if $\xi_1+\xi_2+\xi_3=0$ and $\xi_1\xi_2+ \xi_1\xi_3 + \xi_2 \xi_3 = 0$. Thus, $\beta(\xi_1,\xi_2,\xi_3,\lambda)$ and also $W$ vanish simultaneously with $\sigma(u) = - P(0,u)= \xi_1\xi_2\xi_3 \alpha_0(\xi_1,\xi_2,\xi_3,\lambda)$.

This representation through a local parameter illustrates Riemann vanishing theorem in detail, and allows to examine the case $u\equiv 0$.
\begin{Lemma}\label{L:PsiC34}
Functions $P(\xi,u)$ and $\partial_{u_1} P(\xi,u)$ vanish identically at $u\equiv 0$. Functions $\partial^2_{u_1} P(\xi,u)$ and $\partial_{u_2} P(\xi,u)$ have a zero of multiplicity $3$ at $\xi=0$. Moreover,
\begin{gather*}
 -\partial^2_{u_1} P(\xi,u)|_{u\equiv 0} = \partial_{u_2} P(\xi,u)|_{u\equiv 0} = \psi(\xi).
\end{gather*}
\end{Lemma}
\begin{proof} In fact, the three functions $\partial^2_{u_1} P(\xi,u)|_{u\equiv 0}$, $\partial_{u_2} P(\xi,u)|_{u\equiv 0}$, and $\psi(\xi)$ considered as functions of a path $\gamma$ from infinity to the point $(x(\xi),y(\xi))$ on $\mathcal{V}_{(3,4)}$ satisfy the following functional equation
\begin{gather}
 \psi(x,y;[\gamma + \chi]) = \psi(x,y;[\gamma]) \nonumber\\
 \hphantom{\psi(x,y;[\gamma + \chi]) =}{}\times \exp \left\{ - \left(\oint_{\chi} \rmd r \right)^t
 \left(\int_{\gamma} \rmd u + \frac{1}{2}\oint_{\chi} \rmd u\right) + \pi \rmi \big(\Arf(\phi_\chi) - \Arf(\phi) \big)\right\},\label{PsiFEq}
\end{gather}
where $\rmi^2=-1$, $\chi$ is a cycle, $\Arf$ stands for Arf invariant, $\phi_\chi$ and $\phi$ are Arf functions of the cycle~$\chi$ and the main Arf function of the curve, respectively, for more details see~\cite{BL2005}. The ratio of any two of the functions $\partial^2_{u_1} P(\xi,u)|_{u\equiv 0}$, $\partial_{u_2} P(\xi,u)|_{u\equiv 0}$, and $\psi(\xi)$ is a rational function on~$\mathcal{V}_{(3,4)}$. It is straightforward to verify that such ratio has no poles and so it is a constant.
\end{proof}

On one hand, we have
\begin{gather*}
 \partial^2_{u_1} P(\xi,u) \big|_{u\equiv 0} = \partial^2_{u_1} \sigma (u) \big|_{u = \mathcal{A}(\xi)}.
\end{gather*}
From the series expansion
\begin{gather}
 \sigma\big(t u_1,t^2 u_2,t^5 u_5\big) = \left(u_5 + u_1 u_2^2 - \frac{u_1^5}{20} \right)t^5 -
 \left(\frac{\lambda_2}{6} u_2^2 u_1^3 - \frac{\lambda_2 u_1^7}{2\cdot 6\cdot 7}\right)t^7 \label{C34Sigma}\\
 \qquad{} + \left(\frac{\lambda_2^2}{12} u_1 u_2^4 + \frac{\lambda_2^2}{5!} u_2^2 u_1^5
 - \frac{\lambda_2^2 u_1^9}{4\cdot 6!} \right) t^9
 - \left(\frac{\lambda_5}{6} u_5 u_2 u_1^3 + \frac{\lambda_5}{15} u_2^5 +
 \frac{3\lambda_5}{7!}u_2 u_1^8 \right)t^{10} + O\big(t^{11}\big)\nonumber
\end{gather}
we calculate directly
\begin{gather}\label{C34Psisigma}
 -\log \partial^2_{u_1} P(\xi,u) \big|_{u\equiv 0} = 3\log \xi - \frac{\lambda_5}{12} \xi^5 + O\big(\xi^6\big).
\end{gather}

On the other hand, using notation $c(\lambda)=(c_1,c_2,c_5)^t$, from \eqref{PsiDef} we obtain
\begin{gather}
 \log \psi(\xi) = 3 \log \xi + c_1 \xi + \frac{1}{6} \big(\lambda_2+3c_2 \big) \xi^2
 + \frac{\lambda_2}{6^2} \big(\lambda_2+3 c_2 \big) \xi^4\nonumber\\
 \hphantom{\log \psi(\xi) =}{} - \frac{1}{5\cdot 6^2} \big(21\lambda_5 +4 c_1 \lambda_2^2 + 36 c_5\big)\xi^5 + O\big(\xi^6\big).\label{C34Psi}
\end{gather}
By Lemma \ref{L:PsiC34} the series \eqref{C34Psisigma} and \eqref{C34Psi} are equal. Thus, we come to a~system of \emph{linear} equations, namely
\begin{gather*}
 c_1 = 0,\qquad \lambda_2+3c_2 =0,\qquad \frac{1}{5\cdot 6^2} \big(21\lambda_5 + 4 c_1 \lambda_2^2 + 36 c_5\big)= \frac{\lambda_5}{12},
\end{gather*}
and find \eqref{C34const}. {This solution is unique.}
\end{proof}

\begin{Remark}\label{R:SigmaDer} The primitive function $\psi(\xi)$ coincides, up to a factor which is a rational number, with the first non-vanishing derivative of sigma function on the Abel's image of $\xi$. The first non-vanishing derivative of sigma function has Sato weight equal to $-g$, the corresponding differential operator has Sato weight $-\wgt \sigma (u)-g$. Recall that Sato weight of sigma function is negative and calculated by the formula $-(n^2-1)(s^2-1)/24$ for $(n,s)$-curve, see~\cite{BEL1999}.
\end{Remark}
\begin{Remark}\label{R:SigmaExpansion}
Regarding series expansion \eqref{C34Sigma} for the sigma function related to $(3,4)$-curve of the form \eqref{C34eq}, it was computed on the base of the theory of multivariate sigma functions presented in \cite{BL2004,BL2008}. In \cite[ref.~15]{EEMOP2007} the reader can find a reference to page \url{http://www.ma.hw.ac.uk/Weierstrass/}, where an expansion for the sigma function related to $(3,4)$-curve with some extra parameters is presented, the expansion is obtained by J.C.~Eilbeck. The case of cyclic $(3,4)$-curve, when $\lambda_2$, $\lambda_5$ and $\lambda_8$ vanish, is considered in~\cite{EEMOP2007}, %an earlier preprint arXiv:0610019v1[math.AG],
where a series expansion for the cyclic $(3,4)$-sigma function is proposed. Taking into account the difference in the number of parameters and signs between equations of $(3,4)$-curves, we compared~\eqref{C34Sigma}, and also higher terms which are not presented here, with the two mentioned expansions, and found that they coincide.
\end{Remark}

\begin{proof}[Proof 2]
Let $u$ and $v$ be the Abel's map images of two non-special divisors on $\mathcal{V}_{(3,4)}$, namely,
\begin{gather}\label{uDefC34}
 u = \sum_{i=1}^g \int_{\infty}^{(x_i,y_i)} \rmd u \qquad \text{and}\qquad v = \sum_{i=1}^g \int_{\infty}^{(z_i,w_i)} \rmd u,
\end{gather}
at that $f(x_i,y_i)=0$ and $f(z_i,w_i)=0$, $i=1,2,3$, and $\sigma(u) \neq 0$, $\sigma(v) \neq 0$. By the residue theorem we have
\begin{gather}\label{ZetaDef}
 \sum_{i=1}^{g} \int^{(x_i,y_i)}_{(z_i,w_i)} \rmd r = -\res_{\xi=0}\left( \mathcal{A}^{\ast}(\xi) \frac{\rmd}{\rmd \xi}\log \frac{P(\xi,u)}{P(\xi,v)}\right).
\end{gather}
Direct calculation using \eqref{C34D1k} and \eqref{IntReg} gives
\begin{gather}\label{zetaC34}
 \sum_{i=1}^{g=3} \int_{(z_i,w_i)}^{(x_i,y_i)} \rmd r = \begin{pmatrix}
 -\zeta_1(u) \\
 -\zeta_2(u) + \wp_{1,1}(u)\\
 -\zeta_5(u) + \mathcal{P}_5(u) \end{pmatrix}
 - \begin{pmatrix}
 -\zeta_1(v) \\
 -\zeta_2(v) + \wp_{1,1}(v)\\
 -\zeta_5(v) + \mathcal{P}_5(v) \end{pmatrix},
\end{gather}
where
\begin{gather*}
\mathcal{P}_5(u) = -\tfrac{5}{12} \lambda_2 \wp _{1,2}(u) -\tfrac{5}{8} \wp_{1,2,2}(u)
-\tfrac{5}{12} \wp_{1,1,1,2}(u) -\tfrac{1}{24} \wp_{1,1,1,1,1}(u)
\\
\hphantom{\mathcal{P}_5(u)}{} = -\tfrac{1}{2} \wp_{1,1}(u)\wp_{1,1,1}(u) - \tfrac{5}{2} \wp_{1,1}(u)\wp_{1,2}(u)
-\tfrac{1}{2} \wp_{1,2,2}(u) + \tfrac{1}{6} \lambda_2 \wp_{1,1,1}(u) - \tfrac{5}{12}\lambda_5.
\end{gather*}
Here we apply the relation \eqref{wp1111RelC34}, and use the following notation
\begin{gather}
\zeta_{i}(u) = \partial_{u_i} \log \sigma(u),\qquad \wp_{i,j}(u) = -\partial_{u_i} \partial_{u_j} \log \sigma(u),\nonumber \\
\wp_{i,j,\dots,k}(u) = - \partial_{u_i} \partial_{u_j} \cdots \partial_{u_k} \log \sigma(u).\label{wpNot}
\end{gather}

\begin{Lemma}\label{L:ind4WPC34}
The following relations for Abelian functions on Jacobian of $\mathcal{V}_{(3,4)}$ hold
\begin{subequations}\label{C34WP4rel}
\begin{align}
 &\wp_{1,1,1,1}(u) = 6 \wp_{1,1}^2(u)-3\wp_{2,2}(u) - 4 \lambda_2 \wp_{1,1}(u),\label{wp1111RelC34} \\
 &\wp_{1,1,1,2}(u) = 6\wp_{1,1}(u)\wp_{1,2}(u) - \lambda_2 \wp_{1,2}(u) + \lambda_5. \label{wp1112RelC34}
\end{align}
\end{subequations}
\end{Lemma}
\begin{Remark} The complete list of relations between Abelian functions on Jacobian of genus~$3$ trigonal curve are obtained in \cite{Eilb2011}, the curve is defined by an equation slightly different from~\eqref{C34eq}.
\end{Remark}
\begin{proof}Differentiating \eqref{zetaC34} over $x_1$ we obtain
\begin{gather*}
 \frac{\rmd r(x_1,y_1,\lambda)}{\rmd x_1} = \partial_{u} \begin{pmatrix}
 -\zeta_1(u) \\
 -\zeta_2(u) + \wp_{1,1}(u)\\
 -\zeta_5(u) + \mathcal{P}_5(u) \end{pmatrix} \frac{\rmd u(x_1,y_1,\lambda)}{\rmd x_1},
\end{gather*}
which produces three relations with rational functions $\mathcal{R}_k$ of order $k$ on $\mathcal{V}_{(3,4)}$
\begin{gather*}
 \mathcal{R}_6(x_1,y_1) = 0, \qquad \mathcal{R}_7(x_1,y_1) = 0, \qquad \mathcal{R}_{10}(x_1,y_1) = 0,
\end{gather*}
where
\begin{subequations}\label{RelC34}
 \begin{gather}
 \mathcal{R}_6(x,y) = x^2 + y \wp_{1,1} + x \wp_{1,2} + \wp_{1,5},\label{R6RelC34}\\
 \mathcal{R}_7(x,y) = 2 x y + y \big(\wp_{1,2} + \wp_{1,1,1}\big) + x \big(\wp_{2,2} + \wp_{1,1,2}\big)
 + \big(\wp_{2,5} + \wp_{1,1,5}\big), \label{R7RelC34}\\
 \mathcal{R}_{10}(x,y) = 5 y x^2 - \tfrac{2}{3} \lambda_2^2 x^2 - \tfrac{2}{3} \lambda_5 \lambda_2 x + \lambda_6 y + y \big(\wp_{1,5} + \partial_{u_1} \mathcal{P}_5\big) \notag\\
\hphantom{\mathcal{R}_{10}(x,y) =}{} + x \big(\wp_{2,5} + \partial_{u_2} \mathcal{P}_5\big) + \big(\wp_{5,5} + \partial_{u_5} \mathcal{P}_5\big).\label{R10RelC34}
 \end{gather}
\end{subequations}
For brevity we omit argument $u$ of Abelian functions. Clearly, the functions \eqref{RelC34} vanish on $(x_2,y_2)$ and $(x_3,y_3)$ as well.

In order to proceed, we introduce a rational function of order $10$ {on $\mathcal{V}_{(3,4)}$}
\begin{gather*}
 \varphi_{10}(x,y) = \mathcal{R}_{10}(x,y) -
 \tfrac{5}{4} \bigl(2x - \wp_{1,2} - \wp_{1,1,1}\bigr) \mathcal{R}_{7}(x,y)\\
 \hphantom{\varphi_{10}(x,y) =}{} + \tfrac{1}{6}\bigl(15 \wp_{2,2} + 15\wp_{1,1,2} + 4\lambda_2^2\bigr) \mathcal{R}_{6}(x,y).
\end{gather*}
In fact, $\varphi_{10}(x,y)=(y,x,1)\alpha(u)$ with a certain vector function $\alpha(u)$. Thus, system $\varphi_{10}(x_1,y_1)= \varphi_{10}(x_2,y_2)= \varphi_{10}(x_3,y_3)=0$ is equivalent to a relation of the form
\begin{gather*}
 \begin{pmatrix}
 y_1 & x_1 & 1 \\ y_2 & x_2 & 1\\ y_3 & x_3 & 1
 \end{pmatrix} \alpha (u) = 0.
\end{gather*}
Suppose $(x_i,y_i)$, $i=1,2,3$, are pairwise distinct, then the determinant $\left|\begin{smallmatrix} y_1 & x_1 & 1 \\ y_2 & x_2 & 1\\ y_3 & x_3 & 1 \end{smallmatrix}\right|$ does not vanish since $\sigma(u)\neq 0$. Therefore, $\alpha(u)=0$. By breaking $\alpha (u)$ into even $\alpha (u) + \alpha (-u)$ and odd $\alpha (u) - \alpha (-u)$ parts we obtain in particular
\begin{subequations}
\begin{gather}
 \mathcal{T}_1 = \wp_{1,1,1,1,1,1} + 15\wp_{1,1,2,2} - 30\wp_{1,1,1}^2 - 24\wp_{1,5} -30 \wp_{1,2}^2\nonumber\\
 \hphantom{\mathcal{T}_1 =}{} - 60\wp_{2,2}\wp_{1,1} - 16\lambda_2^2 \wp_{1,1} - 24 \lambda_6 = 0, \label{EvenRel1C34}\\
 \mathcal{T}_2 = \wp_{1,1,1,1,1,2} + 15\wp_{1,2,2,2} - 30\wp_{1,1,1}\wp_{1,1,2} + 36\wp_{2,5}\nonumber\\
 \hphantom{\mathcal{T}_2 =}{} - 90 \wp_{2,2}\wp_{1,2} - 16\lambda_2^2 \wp_{1,2} + 16 \lambda_2 \lambda_5 = 0 , \label{EvenRel2C34}\\
 \mathcal{T}_3 = \wp_{1,1,1,1,2} - 6\wp_{1,1}\wp_{1,1,2} - 6\wp_{1,2}\wp_{1,1,1} + \lambda_2 \wp_{1,1,2} = 0. \label{OddRelC34}
\end{gather}
\end{subequations}

From \eqref{R6RelC34} we see that $\phi_{6}(x,y;u) = x^2 + \wp_{1,1}(u) y + \wp_{1,2}(u) x + \wp_{1,5}(u)$ is even in $u$
and has $2g=6$ roots $(x_i,y_i)$, $i=1,\dots,6$, on curve $\mathcal{V}_{(3,4)}$. At $(x_1,y_1)$, $(x_2,y_2)$, and $(x_3,y_3)$ the function
$\phi_{7}(x,y;u) = 2 x y + \big(\wp_{1,2}(u) + \wp_{1,1,1}(u)\big) y + \big(\wp_{2,2}(u) + \wp_{1,1,2}(u)\big) x
 + \big(\wp_{2,5}(u) + \wp_{1,1,5}(u)\big)$, cf.~\eqref{R7RelC34}, vanishes. At the same time, the function $\phi_{7'}(x,y;u) = \phi_{7}(x,y;-u) = 2 x y + \big(\wp_{1,2}(u) - \wp_{1,1,1}(u)\big) y + \big(\wp_{2,2}(u) - \wp_{1,1,2}(u)\big) x + \big(\wp_{2,5}(u) - \wp_{1,1,5}(u)\big)$ vanishes
at $(x_4,y_4)$, $(x_5,y_5)$, and $(x_6,y_6)$. Consequently, the ratio $\phi_{7}(x,y;u)\phi_{7'}(x,y;u)/\phi_{6}(x,y;u)$ has no poles on $\mathcal{V}_{(3,4)}$ except the pole of order $8$ at infinity, which means that a decomposition
\begin{gather*}
 \phi_{7}\phi_{7'} - \big(a_0 y^2 + a_1 x y + a_2 x^2 + a_4 y + a_5 x + a_8\big)\phi_{6} + b_2 f(x,y) = 0
\end{gather*}
exists. Coefficients of monomials $x^i y^j$ yield an overdetermined system of $13$ linear equations with respect to $a_0$, $a_1$, $a_2$, $a_4$, $a_5$, $a_8$, $b_2$. Its compatibility condition includes in particular
\begin{subequations}
\begin{gather}
\mathcal{T}_4 = \wp _{1,1,1}^2 - 4 \wp _{1,1}^3 + 4 \wp _{2,2} \wp _{1,1}
 - \wp _{2,1}^2 + 4 \wp _{5,1} + 4 \lambda _2 \wp _{1,1}^2 = 0, \label{wp111SqRelC34}\\
\mathcal{T}_5 = \wp _{1,1,1} \wp _{2,1,1} - 4 \wp_{2,1} \wp_{1,1}^2 + \wp_{2,1} \wp_{2,2} - 2 \wp_{5,2}
 + 2 \lambda_2 \wp _{2,1} \wp_{1,1} - 2 \lambda_5 \wp _{1,1} = 0. \label{wp111wp112RelC34}
\end{gather}
\end{subequations}

Taking a linear combination of derivatives of \eqref{EvenRel1C34}, \eqref{EvenRel2C34}, and~\eqref{wp111SqRelC34} we come to
\begin{gather*}
 \partial_{u_2} \mathcal{T}_1 - \partial_{u_1} \mathcal{T}_2 + 15 \partial_{u_2} \mathcal{T}_4 = 30 \wp_{1,1,2}
 \bigl(\wp_{1,1,1,1} - 6\wp_{1,1}^2 + 3\wp_{2,2} +4\lambda_2 \wp_{1,1}\bigr) = 0.
\end{gather*}
Since $\wp_{1,1,2}(u)$ is not identically zero, this implies \eqref{wp1111RelC34}.

Next, denote $\mathcal{T}_6 = \wp_{1,1,1,1} - 6\wp_{1,1}^2 + 3\wp_{2,2} +4\lambda_2 \wp_{1,1}$. Relation \eqref{wp1112RelC34} follows from
\begin{gather*}
 \partial_{u_2} \mathcal{T}_1 - \partial_{u_1} \mathcal{T}_2 - 30\partial_{u_1} \mathcal{T}_5 +
 20 \wp_{1,1} \bigl(\partial_{u_2}\mathcal{T}_6 - \mathcal{T}_3\bigr)\\
 \qquad{} = -60 \wp_{1,1,1} \bigl(\wp_{1,1,1,2} - 6\wp_{1,1}\wp_{1,2} + 3\lambda_2 \wp_{1,2} - \lambda_5\bigr) = 0,
\end{gather*}
by similar reason.
\end{proof}

In the case of $\mathcal{V}_{(3,4)}$ we have the equality, for more detail see \cite{BL2005},
\begin{gather}\label{BLRC34}
 \frac{\sigma\big(u-\mathcal{A}(\xi)\big)\sigma\big(u+\mathcal{A}(\xi)\big)} {\psi^2(\xi)\sigma^2(u)} = \phi_{6}\big(x(\xi),y(\xi);u\big),
\end{gather}
where \eqref{xyParamC34} is applied. Indeed, both left and right hand sides are rational functions on the curve, and vanish at $2g=6$ points which are Abel's map pre-images of $u$ and $-u$. Comparing the leading terms of expansions in the vicinity of $\xi=0$ we see that the functions are equal. The expansion gives
\begin{gather*}
 0= \frac{\sigma\big(u-\mathcal{A}(\xi)\big)\sigma\big(u+\mathcal{A}(\xi)\big)}
 {\psi^2(\xi)\sigma^2(u)} - \phi_{6}\big(x(\xi),y(\xi);u\big)= -2 c_1 \xi^{-5}
 + \left(2c_1^2 - c_2 -\frac{\lambda_2}{3}\right)\xi^{-4}
\\ \hphantom{0=}{} + c_1 (\cdots) \xi^{-3}+
 \left(c_1 \big(\cdots \big) + \frac{c_2^2}{2} + \frac{\lambda_2}{6} c_2 + c_2 \wp_{1,1}
 + \frac{1}{2}\wp_{1,1}^2 -\frac{1}{4}\wp_{2,2} -\frac{1}{12}\wp_{1,1,1,1} \right)\xi^{-2}
 \\ \hphantom{0=}{} - \left(c_1 (\cdots ) - \frac{2}{5}c_5
 -\frac{7\lambda_5}{30} - \left(c_2 + \frac{\lambda_2}{6}\right) \wp_{1,2}
 - \wp_{1,1} \wp_{1,2} + \frac{1}{6}\wp_{1,1,1,2}\right)\xi^{-1} + O(1).
\end{gather*}
Applying \eqref{C34WP4rel} we find \eqref{C34const}.
\end{proof}

\begin{Remark}The equality \eqref{BLRC34} produces bilinear Hirota type equations, which give relations between Abelian functions. Note that
only correct choice of the regularization constant guarantees consistency of relations. So the Hirota type equations carry information about the regularization constant.
\end{Remark}
\begin{thebibliography}{99}
\bibitem[1]{BakerAF}
Baker H.F., Abelian functions: Abel's theorem and the allied theory of theta
 functions, \textit{Cambridge Mathematical Library}, Cambridge University Press,
 Cambridge, 1995.

\bibitem[2]{BEL1999}
Buchstaber V.M., Enolskii V.Z., Leykin D.V., Rational analogues of abelian
 functions, \href{https://doi.org/10.1007/BF02465189}{\textit{Funct. Anal. Appl.}} \textbf{33} (1999), 83--94.

\bibitem[4]{BEL2012}
Buchstaber V.M., Enolskii V.Z., Leykin D.V., Multi-dimensional sigma-functions,
 \href{https://arxiv.org/abs/1208.0990}{arXiv:1208.0990}.

\bibitem[5]{BL2004}
Bukhshtaber V.M., Leykin D.V., Heat equations in a nonholomic frame,
 \href{https://doi.org/10.1023/B:FAIA.0000034039.92913.8a}{\textit{Funct. Anal. Appl.}} \textbf{38} (2004), 88--101.

\bibitem[6]{BL2005}
Bukhshtaber V.M., Leykin D.V., Addition laws on {J}acobian varieties of plane
 algebraic curves, \textit{Proc. Steklov Inst. Math.} \textbf{251} (2005),
 49--120.

\bibitem[7]{BL2008}
Bukhshtaber V.M., Leykin D.V., Solution of the problem of differentiation of
 Abelian functions over parameters for families of {$(n,s)$}-curves,
 \href{https://doi.org/10.1007/s10688-008-0040-4}{\textit{Funct. Anal. Appl.}} \textbf{42} (2008), 268--278.

\bibitem[8]{Eilb2011}
Eilbeck J.C., England M., \^Onishi Y., Abelian functions associated with genus
 three algebraic curves, \href{https://doi.org/10.1112/S1461157010000355}{\textit{LMS~J. Comput. Math.}} \textbf{14} (2011),
 291--326, \href{https://arxiv.org/abs/1008.0289}{arXiv:1008.0289}.

\bibitem[9]{EEMOP2007}
Eilbeck J.C., Enolski V.Z., Matsutani S., \^Onishi Y., Previato E., Abelian
 functions for trigonal curves of genus three, \href{https://doi.org/10.1093/imrn/rnm140}{\textit{Int. Math. Res. Not.}}
 \textbf{2008} (2008), Art. ID rnm~140, 38~pages, \href{https://arxiv.org/abs/math.AG/0610019}{math.AG/0610019}.

\bibitem[10]{EEL2000}
Eilbeck J.C., Enolskii V.Z., Leykin D.V., On the {K}leinian construction of
 abelian functions of canonical algebraic curves, in S{IDE} {III}~--
 Symmetries and Integrability of Difference Equations ({S}abaudia, 1998),
 \textit{CRM Proc. Lecture Notes}, Vol.~25, Amer. Math. Soc., Providence, RI,
 2000, 121--138.

\bibitem[11]{Hodge}
Hodge W.V.D., Atiyah M.F., Integrals of the second kind on an algebraic
 variety, \href{https://doi.org/10.2307/2007100}{\textit{Ann. of Math.}} \textbf{62} (1955), 56--91.

\bibitem[12]{N2010}
Nakayashiki A., On algebraic expressions of sigma functions for {$(n,s)$}
 curves, \href{https://doi.org/10.4310/AJM.2010.v14.n2.a2}{\textit{Asian~J. Math.}} \textbf{14} (2010), 175--211,
 \href{https://arxiv.org/abs/0803.2083}{arXiv:0803.2083}.

\bibitem[14]{Whittaker}
Whittaker E.T., Watson G.N., A course of modern analysis, \href{https://doi.org/10.1017/CBO9780511608759}{\textit{Cambridge
 Mathematical Library}}, Cambridge University Press, Cambridge, 1996.

\end{thebibliography}
\end{document}
